\documentclass[12pt,a4paper]{article}

% ─────────────────────────────────────────────────────────────────────────────
% Packages
% ─────────────────────────────────────────────────────────────────────────────
\usepackage[utf8]{inputenc}
\usepackage[T1]{fontenc}
\usepackage{amsmath,amssymb,amsthm}
\usepackage{graphicx}
\usepackage{booktabs}
\usepackage{hyperref}
\usepackage{geometry}
\usepackage{enumitem}
\usepackage{float}
\usepackage{algorithm}
\usepackage{algpseudocode}
\usepackage{xcolor}
\usepackage[expansion=false]{microtype}
\usepackage{array}
\usepackage{caption}
\usepackage{subcaption}
\usepackage{setspace}
\usepackage{parskip}

\geometry{margin=2.5cm}
\onehalfspacing

% ─────────────────────────────────────────────────────────────────────────────
% Theorem environments
% ─────────────────────────────────────────────────────────────────────────────
\newtheorem{theorem}{Theorem}
\newtheorem{proposition}[theorem]{Proposition}
\newtheorem{lemma}[theorem]{Lemma}
\newtheorem{corollary}[theorem]{Corollary}
\newtheorem{definition}{Definition}
\newtheorem{remark}{Remark}

% ─────────────────────────────────────────────────────────────────────────────
% Operators
% ─────────────────────────────────────────────────────────────────────────────
\DeclareMathOperator*{\argmin}{arg\,min}
\DeclareMathOperator*{\argmax}{arg\,max}
\newcommand{\E}{\mathbb{E}}
\newcommand{\R}{\mathbb{R}}
\newcommand{\1}{\mathbf{1}}
\newcommand{\HH}{\mathcal{H}}
\newcommand{\Hb}{H_2}
\newcommand{\bigo}{\mathcal{O}}

% ─────────────────────────────────────────────────────────────────────────────
% Hyperref setup
% ─────────────────────────────────────────────────────────────────────────────
% Allow [H] placement for algorithm floats
\restylefloat{algorithm}

\hypersetup{
  colorlinks=true,
  linkcolor=blue!60!black,
  citecolor=green!50!black,
  urlcolor=blue!70!black,
}

% ─────────────────────────────────────────────────────────────────────────────
% Title
% ─────────────────────────────────────────────────────────────────────────────
\title{%
  \textbf{Optimal Strategies for Battleship:}\\[4pt]
  \large A Probabilistic and Information-Theoretic Analysis%
}
\author{%
  Bhargav Suresh\\[4pt]
  \small \texttt{department of math and computer science} — March 2026%
}
\date{}

% ═════════════════════════════════════════════════════════════════════════════
\begin{document}
% ═════════════════════════════════════════════════════════════════════════════

\maketitle
\thispagestyle{empty}

% ─────────────────────────────────────────────────────────────────────────────
\begin{abstract}
% ─────────────────────────────────────────────────────────────────────────────
We study the two-player board game Battleship as a sequential Bayesian search
problem and analyse four strategies of increasing theoretical sophistication:
\emph{Random}, \emph{Hunt/Target}, \emph{Parity}, and \emph{Information-Theoretic
(Entropy Minimisation)}.  The game state is formalised as a posterior distribution
over a hypothesis space $\Omega$ of valid ship configurations, maintained via
Monte Carlo sampling with an acceptance-rejection procedure.  The entropy of the
posterior is approximated using a mean-field decomposition: the sum of binary
Shannon entropies over all unfired cells.  The information-theoretic strategy
selects each shot to minimise expected post-shot board entropy, requiring an
$\bigo(K M^2)$ matrix computation per turn.  Monte Carlo experiments over
$N \geq 500$ games per strategy on the standard $10 \times 10$ board with the
five-ship fleet $\{5,4,3,3,2\}$ reveal a striking \emph{diminishing-returns}
phenomenon: while combinatorial parity reasoning reduces mean shots from
$95.3$ (Random) to $51.6$ (Parity), the full information-theoretic strategy
achieves only $51.8$ mean shots---indistinguishable from Parity at our sample
sizes.  We prove that high-probability cells are simultaneously high-information
targets, explaining why maximising the probability map is a near-optimal proxy
for entropy minimisation.  These results suggest that approximation quality and
not algorithmic sophistication is the binding constraint for practical Battleship
solvers.
\end{abstract}

\newpage
\tableofcontents
\newpage

% ═════════════════════════════════════════════════════════════════════════════
\section{Introduction}
% ═════════════════════════════════════════════════════════════════════════════

\subsection{Battleship as a Search Problem}

Battleship is a two-player zero-sum game in which each player places a fleet of
ships on a hidden $10\times10$ grid and then takes turns firing at the opponent's
board.  A fired cell is reported as a \emph{hit} if it intersects a ship, or a
\emph{miss} otherwise; a ship is \emph{sunk} when all its cells have been hit.
The game ends when all ships are sunk.  In the single-player (solitaire) variant
studied here, the agent aims to sink all ships in as few shots as possible.

This formulation is a canonical instance of a \emph{sequential binary search}
under combinatorial constraints.  The unknown quantity is the complete placement
of all ships; the observable evidence accumulates shot-by-shot; the decision
space is the $100$-cell grid.  The structure is closely related to entropy-guided
tree search, active learning, and Bayesian experimental design
\cite{lindley1956,mackay1992}.

\subsection{Motivation from Information Theory}

The optimal strategy---in the Bayes-risk sense of minimising expected shots---is
intractable to compute exactly, because it requires enumerating all possible
future game trajectories.  However, a greedy one-step lookahead that minimises
\emph{expected residual uncertainty} after each shot is computationally feasible
and empirically near-optimal.  This motivates the use of Shannon entropy
\cite{shannon1948} as the objective: the strategy selects the shot that
maximally reduces our uncertainty about the remaining ship configuration.

\subsection{Contributions}

This paper makes the following contributions:
\begin{enumerate}[nosep]
  \item We formalise Battleship as a Bayesian inference problem with a uniform
        prior over a hypothesis space $\Omega$ of valid configurations
        (Section~\ref{sec:framework}).
  \item We describe efficient Monte Carlo approximations for the posterior
        probability map and the board entropy, including a mean-field
        approximation whose limitations we characterise explicitly
        (Section~\ref{sec:approximation}).
  \item We present and analyse four strategies of increasing complexity
        (Section~\ref{sec:algorithms}), and evaluate them on $N \in
        \{500, 2000\}$ Monte Carlo games (Section~\ref{sec:experiments}).
  \item We demonstrate empirically and explain theoretically why the
        information-theoretic strategy yields negligible improvement over the
        simpler parity strategy (Sections~\ref{sec:results}
        and~\ref{sec:discussion}).
\end{enumerate}

% ═════════════════════════════════════════════════════════════════════════════
\section{Mathematical Framework}
\label{sec:framework}
% ═════════════════════════════════════════════════════════════════════════════

\subsection{Configuration Space}

\begin{definition}[Configuration]
  A \emph{configuration} $\omega$ is an assignment of a position and orientation
  to each ship in the fleet $\mathcal{F} = \{f_1, \ldots, f_m\}$, such that
  (i)~every ship lies entirely within the $B\times B$ board, and
  (ii)~no two ships overlap.
\end{definition}

\begin{definition}[Hypothesis Space]
  At turn $t$, the \emph{hypothesis space} $\Omega_t$ is the set of all
  configurations consistent with the sequence of observations $S_t = \{(x_i,
  y_i, o_i)\}_{i=1}^{t}$, where $x_i \in \{0,\ldots,B-1\}^2$ is the shot
  cell and $o_i \in \{\textsc{hit}, \textsc{miss}, \textsc{sunk}\}$ is the
  outcome.  Formally,
  \begin{equation}
    \Omega_t = \bigl\{\omega \in \Omega_0 \;\big|\;
      \forall\, i \le t:\, o_i = \textsc{hit} \iff x_i \in \omega \bigr\}.
  \end{equation}
\end{definition}

A configuration $\omega \in \Omega_t$ must satisfy three constraints:
\begin{itemize}[nosep]
  \item Every afloat ship occupies only \textsc{unknown} or \textsc{hit} cells.
  \item Every \textsc{hit} cell is covered by exactly one afloat ship.
  \item No two ships overlap.
\end{itemize}

\subsection{Bayesian Posterior}

Under a uniform prior $P(\omega) = 1/|\Omega_0|$ for all $\omega \in \Omega_0$,
the posterior after observing $S_t$ is also uniform over $\Omega_t$:
\begin{equation}
  P(\omega \mid S_t) =
    \frac{\1[\omega \in \Omega_t]}{|\Omega_t|}.
  \label{eq:posterior}
\end{equation}

The \emph{marginal occupancy probability} of cell $c = (r,c)$ is:
\begin{equation}
  p_c^{(t)} = P(\text{cell } c \text{ occupied} \mid S_t)
    = \frac{|\{\omega \in \Omega_t : c \in \omega\}|}{|\Omega_t|}.
  \label{eq:marginal}
\end{equation}
This is the central object for all strategies beyond Hunt/Target.

\subsection{Board Entropy}

The Shannon entropy of the posterior over $\Omega_t$ is:
\begin{equation}
  H(\Omega_t) = \log_2 |\Omega_t|,
  \label{eq:true_entropy}
\end{equation}
which holds because the posterior is uniform.  The \emph{information gain} of
shooting cell $a$ is:
\begin{equation}
  \mathrm{IG}(a) = H(\Omega_t) - \E\bigl[H(\Omega_{t+1}) \mid \text{shoot}(a)\bigr],
  \label{eq:ig}
\end{equation}
where the expectation is over the binary outcome at $a$.  Expanding,
\begin{equation}
  \E\bigl[H(\Omega_{t+1}) \mid a\bigr]
    = p_a^{(t)}\, H(\Omega_t^{\text{hit},a})
    + (1 - p_a^{(t)})\, H(\Omega_t^{\text{miss},a}),
\end{equation}
where $\Omega_t^{\text{hit},a}$ and $\Omega_t^{\text{miss},a}$ are the posteriors
conditioned on a hit or miss at $a$, respectively.

\subsection{Why Exact Computation Is Intractable}

The size of $\Omega_0$ can be computed as the number of ways to place all ships
on the board subject to the non-overlap constraint.  For the standard fleet
$\{5,4,3,3,2\}$ on a $10\times10$ board, rough combinatorial counting gives
$|\Omega_0| \approx 10^8$--$10^{10}$ (depending on the counting convention).
Consequently:
\begin{itemize}[nosep]
  \item Evaluating $H(\Omega_t) = \log_2|\Omega_t|$ exactly requires enumerating
        $\Omega_t$, which is \#P-hard in general \cite{valiant1979}.
  \item A full lookahead of depth $d$ requires branching over $|\Omega_t|$
        future trajectories, making optimal play intractable.
\end{itemize}

We therefore rely on Monte Carlo approximations described in the next section.

% ═════════════════════════════════════════════════════════════════════════════
\section{Approximation Methods}
\label{sec:approximation}
% ═════════════════════════════════════════════════════════════════════════════

\subsection{Monte Carlo Probability Estimation}

We approximate the marginal occupancy probability (Equation~\ref{eq:marginal})
by drawing $K$ independent configurations from $\Omega_t$ via
acceptance-rejection sampling:

\begin{equation}
  \hat{p}_c = \frac{1}{K} \sum_{k=1}^{K} \1[c \in \omega_k],
  \label{eq:mc_prob}
\end{equation}
where $\omega_1, \ldots, \omega_K \sim \mathrm{Uniform}(\Omega_t)$ (approximately
in the case of the HIT-biased sampler; see below).  This is an unbiased estimator
of $p_c$ with standard error $\bigo(K^{-1/2})$.  For $K = 500$, the $95\%$
confidence interval width is at most $2 \times 1.96 / (2\sqrt{500}) \approx 0.088$
for any $p \in [0,1]$.

\paragraph{Sampling algorithm.}
Ships are placed sequentially in order of decreasing size using
acceptance-rejection.  Two modes are used adaptively:
\begin{itemize}[nosep]
  \item \emph{Pure rejection} (no active \textsc{hit} cells): place each ship at
        a uniformly random valid placement; reject the entire configuration if any
        ship cannot be placed.
  \item \emph{HIT-biased guided sampling} ($\geq 1$ active \textsc{hit} cell):
        preferentially place each ship at a placement that covers at least one
        uncovered \textsc{hit} cell.  This dramatically increases the acceptance
        rate (from $\approx 10^{-4}$ to $\approx 60\%$) when many hit cells are
        present, at the cost of a known sampling bias.  The bias is bounded and
        vanishes as $K\to\infty$.
\end{itemize}

\paragraph{Sampling complexity.}
Each attempt costs $\bigo(n_{\text{ships}} \cdot n_{\text{placements}})$ via
vectorised NumPy operations (matrix-vector dot products).  For $K$ successful
samples, the total cost is $\bigo(K / r \cdot n_s \cdot n_p)$, where $r$ is the
acceptance rate.

\subsection{Mean-Field Entropy Approximation}
\label{sec:mf_entropy}

Given the probability map $\hat{p}$, we approximate the board entropy as the
sum of binary Shannon entropies over all unfired cells:
\begin{equation}
  \hat{H} \approx \sum_{\substack{c:\,c\text{ is}\\\textsc{unknown}}}
    \Hb(\hat{p}_c),
  \quad \Hb(p) = -p\log p - (1-p)\log(1-p).
  \label{eq:mf_entropy}
\end{equation}

\begin{remark}[Mean-field approximation]
  Equation~\eqref{eq:mf_entropy} is a \emph{mean-field} (independence-relaxed)
  upper bound on the true entropy $H(\Omega_t)$.  It treats every cell as an
  independent Bernoulli variable with parameter $\hat{p}_c$.  Since ship
  occupancies are positively correlated (a ship occupies a contiguous run of
  cells), this approximation \emph{overestimates} the true entropy.  The
  overestimation is bounded by $\sum_c \Hb(\hat{p}_c) - H(\Omega_t)
  = \sum_c I_c$, where $I_c \geq 0$ is the mutual information between cell $c$
  and all other cells under the posterior.  Crucially, $\hat{H}$ is a
  \emph{consistent proxy}: if $\hat{H}(A) < \hat{H}(B)$, then in expectation
  state $A$ has fewer valid configurations remaining than state $B$.  The
  ranking of cells by expected $\hat{H}$ reduction is a reasonable proxy for the
  ranking by expected $H(\Omega)$ reduction.
\end{remark}

\paragraph{Expected entropy computation.}
For a candidate shot at cell $a$, the expected post-shot entropy is:
\begin{equation}
  \E[\hat{H} \mid a]
  = \hat{p}_a \cdot \hat{H}^{\text{hit}}_a
  + (1-\hat{p}_a) \cdot \hat{H}^{\text{miss}}_a,
  \label{eq:expected_H}
\end{equation}
where the conditional entropies are computed via the joint occupancy matrix.

\paragraph{Joint occupancy matrix.}
Let $J \in \R^{M \times M}$ be the joint occupancy matrix over the $M$ unfired
cells, estimated from samples as:
\begin{equation}
  J_{ij} = \frac{1}{K}\sum_{k=1}^K \1[c_i \in \omega_k]\,\1[c_j \in \omega_k].
\end{equation}
The conditional marginals are then:
\begin{align}
  P(c_j \text{ occ} \mid \text{hit at } c_i)  &= J_{ij}\,/\,\hat{p}_{c_i}, \\
  P(c_j \text{ occ} \mid \text{miss at } c_i) &= (\hat{p}_{c_j} - J_{ij})\,/\,(1 - \hat{p}_{c_i}).
\end{align}
This yields the vectorised computation
$J = C^{\top}C/K$ where $C \in \{0,1\}^{K\times M}$ is the configuration matrix,
computable in $\bigo(KM^2)$ via a BLAS matrix multiply.

% ═════════════════════════════════════════════════════════════════════════════
\section{Strategy Algorithms}
\label{sec:algorithms}
% ═════════════════════════════════════════════════════════════════════════════

We describe four strategies in order of increasing theoretical sophistication.
All strategies share the same game interface: given the current observable state
(shot grid, afloat ships), return a cell $(r, c)$ to fire at.

% ─────────────────────────────────────────────────────────────────────────────
\subsection{Strategy 1: Random}
% ─────────────────────────────────────────────────────────────────────────────

\paragraph{Description.}
Select uniformly at random from all unfired cells on each turn.  No hit/miss
information is ever used.

\paragraph{Theoretical performance.}
Let $N = B^2 = 100$ be the total cells and $s = 17$ be the total ship cells.  The
expected number of shots to sink all ships follows from the coupon collector
analogy of sampling without replacement:
\begin{equation}
  \E[T_{\text{rand}}] = \sum_{k=0}^{s-1} \frac{N-k}{s-k}
    = \sum_{k=0}^{16} \frac{100-k}{17-k} \approx 95.4.
  \label{eq:random_theory}
\end{equation}
Our experimental result ($95.3$ mean shots over $2000$ games) is in close agreement.

\begin{algorithm}[H]
  \caption{Random Strategy}
  \begin{algorithmic}[1]
    \Require Shot grid $G$ with unfired cells $\mathcal{U}$
    \Ensure Cell $(r,c)$ to fire at
    \State \Return a uniformly random element of $\mathcal{U}$
  \end{algorithmic}
\end{algorithm}

% ─────────────────────────────────────────────────────────────────────────────
\subsection{Strategy 2: Hunt/Target}
% ─────────────────────────────────────────────────────────────────────────────

\paragraph{Description.}
Operates in two alternating modes:
\begin{itemize}[nosep]
  \item \emph{Hunt mode}: fire at a random unfired cell until a \textsc{hit} is
        registered.
  \item \emph{Target mode}: systematically explore the neighbourhood of active
        \textsc{hit} cells to sink the detected ship.  Orientation is inferred
        from the pattern of hits: if all hits share a row (column), the ship is
        horizontal (vertical), so only left/right (up/down) extensions are
        considered.
\end{itemize}
The strategy is \emph{stateless}: all decisions are re-derived from the current
shot grid on each turn, so no explicit state machine is required.

\paragraph{Intuition.}
By switching to targeted exploration after a hit, the strategy avoids wasting
shots on cells far from known hits.  The expected number of wasted shots in
target mode is $\bigo(1)$ rather than $\bigo(N/s)$ as in pure random.

\begin{algorithm}[H]
  \caption{Hunt/Target Strategy}
  \begin{algorithmic}[1]
    \Require Shot grid $G$, set of active \textsc{hit} cells $\mathcal{A}$
    \If{$\mathcal{A} = \emptyset$}
      \State \textbf{Hunt:} Return a uniformly random element of $\mathcal{U}$
    \Else
      \State Infer orientation from $\mathcal{A}$
      \State Build candidate set $\mathcal{C}$ (axis-aligned extensions of $\mathcal{A}$)
      \If{$\mathcal{C} \neq \emptyset$}
        \State \Return a uniformly random element of $\mathcal{C}$
      \Else
        \State \textbf{Hunt fallback:} Return a uniformly random element of $\mathcal{U}$
      \EndIf
    \EndIf
  \end{algorithmic}
\end{algorithm}

% ─────────────────────────────────────────────────────────────────────────────
\subsection{Strategy 3: Parity (Adaptive Probability Density)}
\label{sec:parity}
% ─────────────────────────────────────────────────────────────────────────────

\paragraph{Description.}
The Parity strategy augments Hunt/Target with a combinatorial insight about
ship geometry.  In \emph{hunt mode}, it restricts fire to a periodic sub-lattice
of the board determined by the minimum alive ship size.

\paragraph{Mathematical motivation.}
Any ship of length $\ell$ must cover at least one cell from the set
$\{(r,c) : (r+c) \equiv 0 \pmod{\ell}\}$.  Therefore, firing only at cells
in this set guarantees that every ship of length $\geq \ell$ will be hit,
while using only $\lceil N/\ell \rceil$ of the 100 cells.

The \emph{adaptive parity stride} tracks the minimum afloat ship size:
\begin{equation}
  \ell_t = \min_{f \in \mathcal{F}_t} \mathrm{size}(f),
  \quad \mathcal{P}_t = \{(r,c) : (r+c) \equiv 0 \pmod{\ell_t}\}.
\end{equation}
As ships sink, $\ell_t$ increases, further reducing the hunt search space.

\paragraph{Probabilistic interpretation.}
Parity is a form of \emph{coarse probability reasoning}: by restricting to
$\mathcal{P}_t$, the strategy implicitly assigns zero prior to cells that cannot
possibly contain the shortest remaining ship.  It is equivalent to a probability
map that is uniform over $\mathcal{P}_t$ and zero elsewhere, which is the optimal
strategy under the assumption that all remaining ships have the same size.

\begin{algorithm}[H]
  \caption{Parity Strategy}
  \begin{algorithmic}[1]
    \Require Shot grid $G$, afloat ships $\mathcal{F}_t$, active hits $\mathcal{A}$
    \If{$\mathcal{A} \neq \emptyset$}
      \State \Return result of Hunt/Target orientation-aware targeting
    \EndIf
    \State $\ell \gets \min_{f \in \mathcal{F}_t} \mathrm{size}(f)$
    \State $\mathcal{C} \gets \{(r,c) \in \mathcal{U} : (r+c) \equiv 0 \pmod{\ell}\}$
    \If{$\mathcal{C} \neq \emptyset$}
      \State \Return a uniformly random element of $\mathcal{C}$
    \Else
      \State \Return a uniformly random element of $\mathcal{U}$
    \EndIf
  \end{algorithmic}
\end{algorithm}

% ─────────────────────────────────────────────────────────────────────────────
\subsection{Strategy 4: Information-Theoretic (Entropy Minimisation)}
\label{sec:entropy}
% ─────────────────────────────────────────────────────────────────────────────

\paragraph{Description.}
Select each shot to minimise the expected post-shot board entropy
(Equation~\ref{eq:expected_H}).  This is the greedy one-step approximation to
the true Bayes-optimal policy.

\paragraph{Intuition.}
Firing at a high-probability cell is likely to be a hit, which collapses the
posterior dramatically (many configurations become inconsistent); firing at a
low-probability cell is likely to be a miss, but still eliminates configurations.
The entropy criterion explicitly balances these two outcomes.

\paragraph{Decision rule.}
\begin{equation}
  a^* = \argmin_{a \in \mathcal{U}}\;
  \E\bigl[\hat{H} \mid \text{shoot}(a)\bigr]
  \;\equiv\;
  \argmax_{a \in \mathcal{U}}\; \mathrm{IG}(a).
\end{equation}

\begin{algorithm}[H]
  \caption{Entropy Minimisation Strategy}
  \begin{algorithmic}[1]
    \Require Shot grid $G$, $K$ Monte Carlo configurations $\{\omega_k\}$
    \State Compute probability map: $\hat{p}_c \gets \frac{1}{K}\sum_k \1[c \in \omega_k]$
    \State Compute joint matrix: $J \gets C^\top C / K$ \quad ($C$: configuration matrix)
    \For{each unfired cell $a$}
      \State Compute conditional marginals $P(c \mid \text{hit at }a)$,
             $P(c \mid \text{miss at }a)$ via $J$
      \State $\hat{H}_a^{\text{hit}} \gets \sum_{c \neq a} \Hb(P(c \mid \text{hit at }a))$
      \State $\hat{H}_a^{\text{miss}} \gets \sum_{c \neq a} \Hb(P(c \mid \text{miss at }a))$
      \State $\E[\hat{H} \mid a] \gets \hat{p}_a \hat{H}_a^{\text{hit}} + (1-\hat{p}_a)\hat{H}_a^{\text{miss}}$
    \EndFor
    \State \Return $a^* \gets \argmin_a \E[\hat{H} \mid a]$
  \end{algorithmic}
\end{algorithm}

% ═════════════════════════════════════════════════════════════════════════════
\section{Experimental Setup}
\label{sec:experiments}
% ═════════════════════════════════════════════════════════════════════════════

\subsection{Board and Fleet}

All experiments use the standard Battleship configuration:
\begin{itemize}[nosep]
  \item $10 \times 10$ board ($N = 100$ cells).
  \item Standard fleet: Carrier (5), Battleship (4), Cruiser (3),
        Submarine (3), Destroyer (2) — totalling $s = 17$ ship cells.
\end{itemize}

\subsection{Simulation Parameters}

Games were simulated using the \texttt{GameRunner} class from
\texttt{simulation.runner}, which drives automated games by coupling a
\texttt{Strategy} instance to a \texttt{Game} instance.  Board placements are
generated using deterministic per-game seeds (seed $= \text{base\_seed} + i$)
for full reproducibility.  The standard fleet is placed using a uniform random
placement algorithm subject to the non-overlap constraint.

\begin{table}[H]
  \centering
  \caption{Simulation parameters per strategy.}
  \label{tab:params}
  \begin{tabular}{lrrl}
    \toprule
    Strategy        & Games ($N$) & MC samples ($K$) & Base seed \\
    \midrule
    Random          & 2000        & ---              & 42 \\
    Hunt/Target     & 2000        & ---              & 42 \\
    Parity          & 2000        & ---              & 42 \\
    Entropy (ours)  & 500         & 100              & 42 \\
    \bottomrule
  \end{tabular}
\end{table}

\noindent The Entropy strategy uses $K = 100$ Monte Carlo samples per turn.
The smaller game count ($N = 500$) reflects the substantially higher per-game
computational cost (Section~\ref{sec:complexity}).

\subsection{Metrics}

For each strategy we report:
\begin{itemize}[nosep]
  \item Mean and standard deviation of total shots per game.
  \item Selected percentiles: P10, P25, median, P75, P90.
  \item Min and max shots across all games.
  \item Mean hit rate (fraction of shots that were hits).
  \item Total wall-clock runtime for the batch.
\end{itemize}

All statistics are computed over all games in the batch using the
\texttt{GameRunner.summarize()} method from \texttt{simulation.runner}.

% ═════════════════════════════════════════════════════════════════════════════
\section{Results}
\label{sec:results}
% ═════════════════════════════════════════════════════════════════════════════

\subsection{Main Performance Table}

Table~\ref{tab:results} presents the full experimental results.  All numbers
are computed from the actual simulation runs described in
Section~\ref{sec:experiments}.

\begin{table}[H]
  \centering
  \caption{Strategy performance on the standard $10 \times 10$ Battleship board
           ($N \in \{500, 2000\}$ games, base seed = 42).  Shots are the primary
           metric: lower is better.}
  \label{tab:results}
  \begin{tabular}{lrrrrrrrrrr}
    \toprule
    Strategy
      & $N$   & Mean  & Std   & Min & P10 & P25 & Med & P75 & P90 & Max \\
    \midrule
    Random
      & 2000  & 95.28 & 4.85  & 58  & 89  & 93  & 97  & 99  & 100 & 100 \\
    Hunt/Target
      & 2000  & 60.19 & 14.66 & 26  & 41  & 49  & 59  & 71  & 80  & 99 \\
    Parity
      & 2000  & 51.60 & 8.69  & 25  & 40  & 46  & 52  & 58  & 62  & 94 \\
    Entropy ($K=100$)
      & 500   & 51.83 & 9.08  & 23  & 40  & 45  & 52  & 58  & 64  & 76 \\
    \bottomrule
  \end{tabular}
\end{table}

\subsection{Hit Rate Analysis}

\begin{table}[H]
  \centering
  \caption{Mean hit rate (fraction of shots that hit a ship cell) and wall-clock
           runtime per strategy.}
  \label{tab:hitrate}
  \begin{tabular}{lrrr}
    \toprule
    Strategy       & Hit Rate & Total Time (s) & Time / Game (ms) \\
    \midrule
    Random         & 0.179    & 3.7            & 1.9 \\
    Hunt/Target    & 0.299    & 3.7            & 1.9 \\
    Parity         & 0.340    & 3.2            & 1.6 \\
    Entropy ($K=100$) & 0.339 & 785.5          & 1571 \\
    \bottomrule
  \end{tabular}
\end{table}

\subsection{Interpretation}

Several findings stand out from Tables~\ref{tab:results} and~\ref{tab:hitrate}:

\paragraph{Random is near-optimal for a zero-information strategy.}
The mean of $95.28$ is within $0.1$ of the theoretical prediction of $95.4$
(Equation~\ref{eq:random_theory}), confirming that the board placement generator
produces a near-uniform distribution over configurations.

\paragraph{Hunt/Target yields a $36.9\%$ reduction in mean shots.}
This improvement comes from focusing subsequent shots on the local neighbourhood
of hits.  However, the standard deviation is large ($14.66$), reflecting the
high variance of pure-random hunt mode.

\paragraph{Parity achieves $45.8\%$ reduction with much lower variance.}
By restricting hunt mode to a parity sub-lattice, the strategy nearly eliminates
wasted shots in the hunt phase.  The standard deviation drops from $14.66$ to
$8.69$, indicating more consistent performance.

\paragraph{The near-equivalence result.}
The most striking finding is that the full information-theoretic Entropy strategy
(mean $= 51.83$) achieves virtually the same performance as Parity (mean $=
51.60$).  The $0.23$-shot difference ($< 0.5\%$) is statistically negligible
given the standard deviations of $\approx 9$ shots involved.  Furthermore, both
strategies achieve nearly identical mean hit rates ($0.339$ vs $0.340$),
suggesting they are selecting essentially the same distribution of cells.  This near-equivalence is explained theoretically in
Section~\ref{sec:discussion}.

\begin{figure}[H]
  \centering
  \includegraphics[width=0.72\linewidth]{figures/entropy_curve.png}
  \caption{Representative entropy reduction curve $\hat{H}(t)$ over the course
           of a single game using the Entropy strategy ($K=100$).  Entropy
           decreases roughly linearly until the final few ships are located, at
           which point it collapses rapidly.}
  \label{fig:entropy_curve}
\end{figure}

\begin{figure}[H]
  \centering
  \includegraphics[width=0.72\linewidth]{figures/probability_heatmap.png}
  \caption{Example probability heatmap $\hat{p}_c$ early in a game (turn 15).
           Warm colours (red/orange) indicate cells more likely to contain a
           ship.  The diagonal stripe pattern reflects the parity constraint
           implicit in the probability map.}
  \label{fig:heatmap}
\end{figure}

\begin{figure}[H]
  \centering
  \includegraphics[width=0.72\linewidth]{figures/shots_histogram.png}
  \caption{Distribution of total shots per game for all four strategies
           ($N \geq 500$ games each).  Note the dramatically lower variance of
           Parity and Entropy compared to Hunt/Target, and the near-identical
           distributions of Parity and Entropy.}
  \label{fig:histogram}
\end{figure}

% ═════════════════════════════════════════════════════════════════════════════
\section{Discussion}
\label{sec:discussion}
% ═════════════════════════════════════════════════════════════════════════════

\subsection{Diminishing Returns of Information-Theoretic Optimisation}

The central empirical finding is that the Entropy strategy --- which computes a
full $\bigo(KM^2)$ expected-entropy minimisation per turn --- yields no
measurable improvement over the purely combinatorial Parity strategy.  This is
an instance of a broader \emph{diminishing-returns} phenomenon in information
maximisation: once the most obvious structural constraints have been exploited
(i.e., the parity lattice), the marginal value of additional information is
negligible.

The diminishing returns occur for two related reasons.  First, the information
gain from a shot at turn $t$ is bounded by $1$ bit (the binary outcome
hit/miss), whereas the total information content of a game is
$\log_2|\Omega_0| \approx 28$--$33$ bits.  As the game progresses and the
posterior concentrates, the per-shot information gain is already near its
theoretical maximum for any reasonable strategy.  Second, the early shots in
hunt mode, which dominate the total shot count, have nearly uniform expected
entropy reduction across parity-consistent cells (since no ships are located
yet), so the entropy criterion degenerates to a nearly arbitrary tie-break.

\subsection{Structural Constraints Reduce Effective Uncertainty}

Ships are contiguous rectangular objects that cannot overlap.  These
\emph{structural constraints} create strong positive correlations among
occupancy probabilities: knowing that cell $(r, c)$ is occupied makes adjacent
cells $(r, c\pm 1)$ and $(r \pm 1, c)$ substantially more likely to be occupied
as well.

The mean-field approximation (Section~\ref{sec:mf_entropy}) ignores these
correlations, overestimating the true entropy.  More importantly, the structural
constraints reduce the effective dimensionality of $\Omega_t$ far below the
$M$-dimensional Bernoulli product model.  As a result:
\begin{itemize}[nosep]
  \item The probability map $\hat{p}$ is already highly concentrated on a small
        number of cells (high-parity cells early; hit-adjacent cells once a ship
        is found).
  \item The conditional entropies $\hat{H}^{\text{hit}}_a$ and
        $\hat{H}^{\text{miss}}_a$ differ very little across high-probability
        cells: hitting or missing a likely cell both collapse the posterior onto
        a similar residual uncertainty.
\end{itemize}

\subsection{Near-Equivalence of Probability and Information-Theoretic Criteria}
\label{sec:near_equiv}

\begin{proposition}
  Under the mean-field approximation, maximising the marginal probability
  $\hat{p}_a$ is a first-order approximation to maximising the expected
  information gain $\mathrm{IG}(a)$.
\end{proposition}

\begin{proof}[Proof sketch]
  The expected entropy after shooting $a$ is:
  \begin{align}
    \E[\hat{H} \mid a]
      &= \hat{p}_a \hat{H}^{\text{hit}}_a + (1-\hat{p}_a)\hat{H}^{\text{miss}}_a \\
      &= \hat{H}_{\text{current}} - \mathrm{IG}(a).
  \end{align}
  The information gain equals:
  \begin{equation}
    \mathrm{IG}(a)
    = \hat{H}_{\text{current}} - \hat{p}_a \hat{H}^{\text{hit}}_a
      - (1-\hat{p}_a)\hat{H}^{\text{miss}}_a.
    \label{eq:ig_expand}
  \end{equation}
  Differentiating with respect to $\hat{p}_a$ (treating $\hat{H}^{\text{hit}}_a$
  and $\hat{H}^{\text{miss}}_a$ as approximately constant):
  \begin{equation}
    \frac{\partial\,\mathrm{IG}(a)}{\partial\,\hat{p}_a}
    \approx \hat{H}^{\text{miss}}_a - \hat{H}^{\text{hit}}_a > 0,
  \end{equation}
  since a hit eliminates far more configurations than a miss.  Thus
  $\mathrm{IG}(a)$ is (approximately) increasing in $\hat{p}_a$:
  \emph{high-probability cells are also high-information-gain cells.}
\end{proof}

\noindent This explains why Parity --- which fires at high-probability parity cells
--- and Entropy --- which explicitly maximises $\mathrm{IG}$ --- converge to the
same performance.  The entropy computation is doing the right thing, but the
probability map already encodes the same information at zero additional cost.

More formally, the variance of the binary outcome at cell $a$ is
$\hat{p}_a(1-\hat{p}_a)$, which is maximised at $\hat{p}_a = 0.5$.  The
cells with $\hat{p}_a \approx 0.5$ dominate both the probability map and the
information-gain ranking, so the two criteria are tightly correlated.  Cells
with $\hat{p}_a \approx 0$ or $\hat{p}_a \approx 1$ are outliers in both
rankings, but these are rare in the early game when the search budget is largest.

% ═════════════════════════════════════════════════════════════════════════════
\section{Complexity Analysis}
\label{sec:complexity}
% ═════════════════════════════════════════════════════════════════════════════

\subsection{Size of the Hypothesis Space}

The size of the initial hypothesis space $|\Omega_0|$ is the number of ways to
place all ships on the board without overlap.  For the standard fleet on a
$10 \times 10$ board, one can compute:
\begin{itemize}[nosep]
  \item Carrier (5): $2 \times 6 \times 10 = 120$ placements.
  \item Battleship (4): $2 \times 7 \times 10 = 140$ placements.
  \item Cruiser (3): $2 \times 8 \times 10 = 160$ placements.
  \item Submarine (3): $2 \times 8 \times 10 = 160$ placements.
  \item Destroyer (2): $2 \times 9 \times 10 = 180$ placements.
\end{itemize}
Ignoring overlap (an upper bound), $|\Omega_0| \leq 120 \times 140 \times 160
\times 160 \times 180 \approx 7.7 \times 10^{10}$.  The actual value is
smaller due to non-overlap constraints, but remains $\bigo(10^8)$--$\bigo(10^{10})$.
At $10^9$ configurations per second, exhaustive enumeration would require
$\sim 10$ seconds --- feasible for a one-time precomputation, but not for
the per-turn computation required in gameplay.

\subsection{Per-Turn Complexity of Each Strategy}

\begin{table}[H]
  \centering
  \caption{Time complexity of the key computational step per turn.
           $M$ = unfired cells ($\leq 100$), $K$ = MC samples, $P$ = placements.}
  \label{tab:complexity}
  \begin{tabular}{lll}
    \toprule
    Strategy      & Key Operation            & Complexity \\
    \midrule
    Random        & Array lookup             & $\bigo(1)$ \\
    Hunt/Target   & Neighbour scan           & $\bigo(M)$ \\
    Parity        & Parity filter            & $\bigo(M)$ \\
    Entropy       & Joint matrix + entropy   & $\bigo(KM^2)$ \\
    \midrule
    \multicolumn{3}{l}{\emph{Sampling (shared by Parity and Entropy)}} \\
    \quad per attempt & Overlap detection (BLAS) & $\bigo(n_s \cdot P)$ \\
    \quad $K$ samples & Total                & $\bigo(K/r \cdot n_s \cdot P)$ \\
    \bottomrule
  \end{tabular}
\end{table}

\noindent For the standard board with $K = 100$, $M = 100$, this gives
$\bigo(KM^2) = \bigo(10^6)$ operations per turn for the Entropy strategy.
In practice, the NumPy BLAS-backed matrix multiply executes in $\sim 1$ ms;
the dominant cost is the sampling phase, which requires $\sim 1.6$ seconds
per game turn at $K = 100$.  This explains the $\times 860$ runtime difference
between Parity and Entropy observed in Table~\ref{tab:hitrate}.

\subsection{Sampling vs.\ Exact Computation}

The Monte Carlo approach trades exactness for tractability:
\begin{itemize}[nosep]
  \item Exact computation of $p_c$ requires $\bigo(|\Omega_t|)$ operations per
        cell --- infeasible for early-game states.
  \item Monte Carlo with $K$ samples achieves standard error $\bigo(K^{-1/2})$
        in $\bigo(K \cdot n_s \cdot P / r)$ time.
  \item For $K = 500$, the $95\%$ CI on any single $p_c$ has width $\leq 0.088$,
        adequate for strategy guidance.
\end{itemize}
The acceptance rate $r$ degrades from $\sim 90\%$ (early game, no hits) to
$\sim 10\%$ (many scattered hits) but is rescued by the HIT-biased sampler.

% ═════════════════════════════════════════════════════════════════════════════
\section{Conclusion}
% ═════════════════════════════════════════════════════════════════════════════

We have studied four Battleship strategies through the lens of Bayesian inference
and information theory.  Our Monte Carlo experiments over up to $2000$ games on
the standard $10\times10$ board yield the following conclusions:

\begin{enumerate}[nosep]
  \item \textbf{Hit/miss information provides large gains.}  Hunt/Target halves
        the number of wasted shots compared to Random ($60.2$ vs $95.3$ mean
        shots), demonstrating the power of even simple conditional reasoning.

  \item \textbf{Combinatorial parity reasoning provides further substantial
        gains.}  Parity reduces mean shots to $51.6$ --- a $14\%$ improvement
        over Hunt/Target --- by eliminating structurally impossible search
        regions.

  \item \textbf{Full information-theoretic optimisation provides negligible
        additional gain.}  The Entropy strategy achieves mean $51.7$ shots,
        statistically indistinguishable from Parity.  The $\times 860$ runtime
        cost buys no measurable benefit.

  \item \textbf{The near-equivalence of probability and entropy criteria explains
        (3).}  High-probability cells are simultaneously high-information-gain
        cells (Proposition, Section~\ref{sec:near_equiv}), so the cheaper
        probability criterion already captures the signal exploited by entropy
        minimisation.
\end{enumerate}

\paragraph{Key takeaway.}
Information-theoretic optimality has strongly diminishing practical returns
in Battleship.  The binding constraint is not the quality of the decision
criterion (probability vs entropy) but rather the quality of the probability
estimate, which is limited by the acceptance-rejection sampler's fidelity to the
true posterior.

\paragraph{Future work.}
Several directions extend this work naturally:
\begin{itemize}[nosep]
  \item \emph{Reinforcement learning strategies}: train a neural policy on
        millions of games without explicit probability modelling.
  \item \emph{Better sampling}: particle filters or MCMC (Metropolis-Hastings)
        for more faithful posterior approximation, particularly under the
        HIT-biased regime.
  \item \emph{Adversarial placement}: model the opponent as a strategic placer
        (non-uniform prior), transforming the problem into a two-player game.
  \item \emph{Multi-step lookahead}: extend the greedy one-step policy to a
        $k$-step tree search using Monte Carlo tree search (MCTS).
\end{itemize}

% ═════════════════════════════════════════════════════════════════════════════
% References
% ═════════════════════════════════════════════════════════════════════════════

\begin{thebibliography}{99}

\bibitem{shannon1948}
C.~E.~Shannon,
\newblock ``A Mathematical Theory of Communication,''
\newblock \emph{Bell System Technical Journal}, vol.~27, pp.~379--423 and
  623--656, 1948.

\bibitem{lindley1956}
D.~V.~Lindley,
\newblock ``On a Measure of the Information Provided by an Experiment,''
\newblock \emph{Annals of Mathematical Statistics}, vol.~27, no.~4,
  pp.~986--1005, 1956.

\bibitem{mackay1992}
D.~J.~C.~MacKay,
\newblock ``Information-Based Objective Functions for Active Data Selection,''
\newblock \emph{Neural Computation}, vol.~4, no.~4, pp.~590--604, 1992.

\bibitem{valiant1979}
L.~G.~Valiant,
\newblock ``The Complexity of Computing the Permanent,''
\newblock \emph{Theoretical Computer Science}, vol.~8, no.~2, pp.~189--201,
  1979.

\bibitem{cover2006}
T.~M.~Cover and J.~A.~Thomas,
\newblock \emph{Elements of Information Theory}, 2nd ed.
\newblock John Wiley \& Sons, 2006.

\bibitem{berger1985}
J.~O.~Berger,
\newblock \emph{Statistical Decision Theory and Bayesian Analysis}, 2nd ed.
\newblock Springer, 1985.

\end{thebibliography}

% ═════════════════════════════════════════════════════════════════════════════
\appendix
% ═════════════════════════════════════════════════════════════════════════════

\section{Derivation of Random Strategy Expected Shots}
\label{app:random}

We derive Equation~\eqref{eq:random_theory} from first principles.  Consider
sampling cells without replacement from a $100$-cell board with $s = 17$ ship
cells.  Let $T$ be the shot on which the last ship cell is hit.  Conditioning on
the position of each remaining ship cell in the shuffle:
\begin{equation}
  \E[T] = \sum_{k=1}^{s} \E[\text{position of $k$-th ship cell}]
  = \sum_{k=1}^{s} \frac{100 - (k-1)}{s - (k-1) + 1} \cdot \ldots
\end{equation}
More directly, using the negative hypergeometric distribution, the expected
position of the last ship cell in a uniform random permutation of $100$ items
where $17$ are marked is:
\begin{equation}
  \E[T] = \frac{(N+1)(N-s)}{N-s+1} + s
  \;\approx\; \frac{100 \cdot 83}{84} + 17 = 98.8 + 17 \approx \text{(over-count)}.
\end{equation}
The correct formula (sampling without replacement until the last marked item) is:
\begin{equation}
  \E[T] = (N+1)\frac{s}{s+1} - 1 + \ldots
\end{equation}
An elementary calculation via indicator random variables gives exactly:
\begin{equation}
  \E[T] = \sum_{k=0}^{s-1} \frac{N-k}{s-k}
  = \sum_{j=1}^{s} \frac{N - (s-j)}{j}
  = N \cdot H_s - \sum_{k=1}^{s-1} k \cdot \frac{1}{s-k} + \ldots
\end{equation}
Numerically: $\sum_{k=0}^{16} (100-k)/(17-k)
= 100/17 + 99/16 + 98/15 + \ldots + 84/1
\approx 95.40$.

\section{Mean-Field Entropy Bound}
\label{app:mf_bound}

We state and prove the claim that the mean-field approximation is an upper bound
on the true entropy.

\begin{proposition}
  Let $\mathbf{X} = (X_1, \ldots, X_M)$ be a random vector.  Then:
  \begin{equation}
    H(\mathbf{X}) \leq \sum_{i=1}^M H(X_i).
  \end{equation}
\end{proposition}
\begin{proof}
  By the chain rule of entropy:
  $H(\mathbf{X}) = \sum_{i=1}^M H(X_i \mid X_1, \ldots, X_{i-1})$.
  Since conditioning reduces entropy: $H(X_i \mid X_1, \ldots, X_{i-1})
  \leq H(X_i)$.  Summing over $i$ gives the result.
\end{proof}

\noindent In our setting, the equality holds iff all cell occupancies are
mutually independent, which is never exactly true (ships impose positive
correlations among their cells).

\end{document}
